% F1-VLA training objective — explanation for a report/thesis.
% Requires: \usepackage{amsmath,amssymb}

\subsection{The F1 training objective}

F1 is trained with two objectives optimised jointly: an \emph{autoregressive
next-scale prediction} loss that teaches the generation expert to imagine what
the scene will look like next, and a \emph{flow-matching} loss that teaches the
action expert to produce motor commands. Intuitively, the model is asked to
answer two questions at once: \emph{``what will I see?''} and \emph{``what
should I do?''}.

\paragraph{Visual foresight as next-token prediction.}
A frozen VQ-VAE encodes each frame into a sequence of discrete visual tokens
drawn from a codebook of $K=4096$ entries, so an image becomes a short
``sentence'' of visual words. Given the history of observations
$\{o_{t-m},\dots,o_t\}$ and the language instruction $l$, the model predicts the
tokens $z_{1:N}$ of the future frame one at a time, each conditioned on the
tokens it has already produced:
\begin{equation}
    \mathcal{L}^{\mathrm{pred}}_{\mathrm{gen}}
    = -\,\mathbb{E}_{\{o_i,l\}\sim\mathcal{D}}
      \left[\sum_{j=1}^{N} \log p_\theta\!\left(z_j \mid \hat{z}_{1:j-1},
      \{o_i\}_{i=t-m}^{t}, l\right)\right].
    \label{eq:gen}
\end{equation}
This is ordinary cross-entropy, i.e.\ ``how surprised is the model by the
correct visual word''. The important design choice is that the conditioning uses
the model's \emph{own} previous predictions $\hat{z}_{1:j-1}$ rather than the
ground-truth ones. Under teacher forcing the model would only ever be asked to
continue perfect prefixes, and would then drift at inference time, when it must
consume its own (possibly wrong) outputs. Conditioning on its own samples during
training removes that train/inference mismatch and is what the authors refer to
as distributional consistency, which in turn stabilises long-horizon rollouts.

\paragraph{Action prediction by flow matching.}
Rather than regressing the action directly, the action expert learns to
transport Gaussian noise into an expert action. For a ground-truth action $a_t$,
a noise sample $\varepsilon\sim\mathcal{N}(0,I)$ and a random interpolation
coefficient $\tau\sim\mathcal{U}(0,1)$, define the intermediate point
$a_t^{\tau} = (1-\tau)\varepsilon + \tau a_t$, which lies on the straight path
from noise to action. The policy $\pi_\theta$ is placed at that point and asked
in which direction to move in order to reach the action; the correct answer is
the constant velocity $a_t - \varepsilon$:
\begin{equation}
    \mathcal{L}_{\mathrm{action}}
    = \mathbb{E}_{\{a_t,o_i,q_t,l\}\sim\mathcal{D}}
      \left\|\pi_\theta\!\left(l, \{o_i\}_{i=t-m}^{t}, q_t, a_t^{\tau}\right)
      - (a_t - \varepsilon)\right\|^2 ,
    \label{eq:action}
\end{equation}
where $q_t$ is proprioception at time $t$. At inference the action is obtained
by sampling noise and integrating the learned vector field over a small number
of Euler steps. The practical benefit over direct regression is multimodality:
when several distinct actions are reasonable in a given state, a regressor is
driven towards their (possibly invalid) mean, whereas a flow model can reach any
of them depending on the sampled noise. The corresponding cost is that rollouts
become stochastic, so evaluation must be averaged over seeds.

\paragraph{Combined objective.}
The two terms are optimised as a weighted sum,
\begin{equation}
    \mathcal{L}_{\mathrm{total}}
    = \mathcal{L}^{\mathrm{pred}}_{\mathrm{gen}}
      + \lambda \, \mathcal{L}_{\mathrm{action}},
    \label{eq:total}
\end{equation}
so that foresight and control are learned in a single representation rather than
in separate stages: the generation term forces the backbone to encode where the
scene is heading, while the action term keeps that representation grounded in
control.

\paragraph{Remark on the implementation.}
In the released code the objective is assembled with the weight placed on the
\emph{generation} term instead of the action term,
\begin{equation*}
    \mathcal{L}_{\mathrm{total}}
    = \mathcal{L}_{\mathrm{action}}
      + \rho \, \mathcal{L}^{\mathrm{pred}}_{\mathrm{gen}},
    \qquad \rho = 0.1 ,
\end{equation*}
which is Eq.~\eqref{eq:total} up to a global factor: dividing by $\rho$ recovers
the paper's form with $\lambda = 1/\rho = 10$. Since a constant rescaling of the
loss only rescales the gradients (and can be absorbed into the learning rate),
the two parameterisations optimise the same objective; the action term is
weighted an order of magnitude more heavily than visual foresight, i.e.\ the
generation loss acts as an auxiliary signal rather than a co-equal one.
Empirically the two terms have very different magnitudes as well
($\mathcal{L}_{\mathrm{action}}\approx 0.2$ versus
$\mathcal{L}^{\mathrm{pred}}_{\mathrm{gen}}\approx 3.7$ in our BridgeData~V2
finetuning run), so after weighting they contribute comparably to the total.
